\documentclass[10pt,a4paper]{article}
\usepackage[margin=25mm]{geometry}
\usepackage{amsmath,amssymb,lmodern}
\setlength{\emergencystretch}{2em}
\begin{document}
\begin{center}
{\Large A definite pencil beyond quotients of individual spectra}\\[4pt]
{\large Counterexample to Conjecture 00000007177}\\[6pt]
ziangni-sys \qquad 4 October 2026
\end{center}

\section*{Precise assertion and matrices}
The source states that generalized eigenvalues have the scale of
the quotient of the two spectra. We refute the individual-spectral
quotient assertion: a pencil eigenvalue need not equal $\alpha/\beta$
for any individual eigenvalues $\alpha$ of $A$ and $\beta\ne0$ of $B$.
Consider
\[
A=\begin{pmatrix}1&0\\0&2\end{pmatrix},\qquad
B=\begin{pmatrix}2&1\\1&3\end{pmatrix},\qquad
v=\begin{pmatrix}1\\-1\end{pmatrix}.
\]
Both matrices are real symmetric positive definite. Indeed, for a
nonzero $z=(x,y)^T$,
\[
z^TAz=x^2+2y^2>0,\qquad
z^TBz=2x^2+2xy+3y^2=x^2+(x+y)^2+2y^2>0.
\]
Thus the example lies within the definite-pair setting.

\section*{An actual pencil eigenvalue}
The standard generalized eigenvalue condition for the pencil
$A-\lambda B$ is a nonzero vector with $Av=\lambda Bv$. Here
\[
Av=\begin{pmatrix}1\\-2\end{pmatrix}=Bv,
\]
so $\lambda=1$ is a genuine generalized eigenvalue.
Neither the vector equation nor definiteness is assumed.

\section*{No quotient of the individual spectra}
For $Az=\alpha z$ with nonzero $z=(x,y)^T$, the two coordinate equations
are $(\alpha-1)x=0$ and $(\alpha-2)y=0$.
At least one coordinate is nonzero, so $\alpha\in\{1,2\}$.
Conversely, the two coordinate vectors realize $1$ and $2$, proving
that the complete individual spectrum of $A$ is $\{1,2\}$.

Neither number is an eigenvalue of $B$. If $Bz=z$, its equations give
$x+y=0$ and $x+2y=0$, forcing $x=y=0$.
If $Bz=2z$, its equations give $y=0$ and $x+y=0$, again forcing zero.
Thus neither equation has a nonzero eigenvector.
Now if $1=\alpha/\beta$ with $\beta\ne0$, then $\alpha=\beta$.
An individual spectral quotient equal to $1$ would therefore require
$A$ and $B$ to share an eigenvalue, which the preceding arguments exclude.
The generalized eigenvalue $1$ is absent from every possible quotient
of their individual spectra, irrespective of how those spectra are ordered.

\section*{Distinction from the Rayleigh quotient}
The generalized Rayleigh quotient uses the same vector in both forms:
\[
R(z)=\frac{z^TAz}{z^TBz}.
\]
It is not a quotient of two individual eigenvalues. In this example
$v^TAv=v^TBv=3$, so $R(v)=1$, exactly as expected.
The counterexample does not dispute this valid variational formulation
or spectral interval bounds. It disproves the source's asserted
individual-spectra quotient description; the separate extremum clause
is unnecessary for refuting the conjunction.

\section*{Formal verification}
Lean defines actual matrices, matrix-vector multiplication, eigenvalues
through nonzero eigenvectors, and the pencil condition $Az=\lambda Bz$.
It proves symmetry and positive definiteness for all nonzero vectors,
the concrete pencil witness, the full individual spectrum of $A$,
the two exclusions for $B$, and impossibility of every quotient
representing $1$. The final theorem packages these properties
without assumed spectral tables or supplied definiteness data.
\end{document}
